\SourceLeaf{351}
In meiner (der ersten) Konstruktion, in der $\delta$ unbestimmt ist, ist die Clise dieselbe wie die von
\[
(a-b)(a-c)(b-c).
\]
Nennt man $e$ den vierten Schnitt[punkt] der Ellipse mit dem Kreis durch $b$, $c$, $d$ und setzt
\[
e = \varepsilon m + \frac{n}{\varepsilon},
\]
wodurch
\[
\beta\gamma\delta\varepsilon = 1
\]
gilt, so ist in \textsc{Newton}s Konstruktion die Clise der Differenz zweier Werte von $t$ dieselbe wie die von
\[
(c-b)(b-e)(c-e).
\]
\Needspace{10\baselineskip}Die obige Formel ($a=\varepsilon\rho$) lässt sich auch besser so schreiben:
\[
t = -\frac{(m-\delta\beta n)(m-\delta\gamma n)(m-\delta\varepsilon n)\rho}{\delta(1-\rho)mn} + \frac{(\beta+\gamma)mm-\dfrac{1}{\varepsilon}mn+\delta nn}{m}
\]
\[
= -\frac{(b-c)(b-e)(c-e)}{\delta(\beta-\gamma)(\beta-\varepsilon)(\gamma-\varepsilon)}\cdot\frac{\rho}{(1-\rho)mn} + \frac{(\beta+\gamma)mm-\dfrac{1}{\varepsilon}mn+\delta nn}{m}.
\]

\Needspace{10\baselineskip}\begin{center}\textbf{BEMERKUNGEN.}\end{center}

Die Notiz [III], die auf einem einzelnen Blatt steht, gibt den Inhalt von Untersuchungen wieder, über die \textsc{Gauss} in dem auf S.\,295 bis 296 dieses Bandes abgedruckten Brief an \textsc{Schumacher} vom 12. Mai 1843 berichtet hat; zur Erläuterung des Textes sei auf diesen Brief verwiesen. Die Notizen [I] und [II], die mit [III] in engem Zusammenhang stehen, finden sich in einem Handbuch und auf einem Zettel; [I] stammt wahrscheinlich schon aus dem Jahr 1831, [II] wohl aus der Zeit um 1840.
\hfill \textsc{Stäckel}.
